% Proof-bearing workbench module.  The standalone root supplies its preamble.
\providecommand{\ZZ}{\mathbb Z}
\providecommand{\ZZpos}{\mathbb Z_{>0}}
\providecommand{\ES}{\operatorname{ES}}
\providecommand{\Rel}{\operatorname{Rel}}
\providecommand{\coker}{\operatorname{coker}}
\providecommand{\eps}{\varepsilon}

\section{The exact marked Erd\H{o}s--Straus--\texorpdfstring{\(S^6\)}{S6} bridge}
\label{sec:esksb-bridge}

\subsection*{Result and scope}
For each prime \(p\equiv1\pmod {12}\), the complete first-half
Erd\H{o}s--Straus divisor selector is in explicit bijection with a
normalized, positively oriented, divisor-supported, and shell-marked slice of
the terminal parameter lattice occurring in the candidate complex structure
on \(S^6\).  The coordinate change is integral in both directions.  More
strongly, the selector relation matrix and the S6 terminal relation matrix
have exactly the same row lattice: they differ by a displayed element of
\(\operatorname{SL}_2(\ZZ)\).  Composing the bijection with the shell
reconstruction gives positive denominators solving \(4/p=1/x+1/y+1/z\), and
every positive solution is detected by at least one of the two tagged slices.

The nonlinear condition \(u\mid a_j^2\), together with the external labels
\((p,j,\eps)\), is indispensable.  If it is removed, each normalized S6 fibre
is an infinite affine line and carries no Erd\H{o}s--Straus existence
information.  A particularly natural unmarked S6 point produces precisely
the already-solved congruence class \(p\equiv13\pmod {24}\) and produces no
selector state when \(p\equiv1\pmod {24}\).  Thus the bridge is an exact
arithmetic equivalence of marked presentations, not a proof of universal
selector nonemptiness and not a geometric morphism from the reciprocal cubic
to the complex threefold.

\subsection{All objects and conventions}

Fix throughout a prime
\[
 p\equiv1\pmod {12},\qquad
 h=\frac{p-1}{12}\in\ZZpos,
\]
and write
\[
 L=\frac{p-1}{4}=3h,
 \qquad J_p=\{j\in\ZZ:0\leq j<L\}.
\]
For \(j\in J_p\), retain the four integers
\begin{equation}
 R_j=4j+3,
 \qquad a_j=3h+j+1=\frac{p+R_j}{4},
 \qquad N_j=pa_j,
 \qquad 4a_j=p+R_j.                         \tag{1.1}
\end{equation}
No label in \((p,j)\) will be suppressed.  The inequalities
\[
 3\leq R_j\leq p-2,
 \qquad \frac{p+3}{4}\leq a_j\leq\frac{p-1}{2}<p
\]
give
\begin{equation}
 \gcd(p,a_j)=\gcd(p,R_j)=\gcd(a_j,R_j)=1.  \tag{1.2}
\end{equation}
Indeed \(p\) is prime and \(0<a_j<p\); then \(R_j=4a_j-p\)
gives the other two equalities.

Let
\[
 \mathscr S(p)=
 \left\{(x,y,z)\in\ZZpos^3:
 4xyz=p(xy+xz+yz)\right\}.
\]
The equation is symmetric, but whenever completeness is invoked we choose the
ordered representative \(x\leq y\leq z\).
Write \(\ES(p)\) for the assertion \(\mathscr S(p)\neq\varnothing\).

For the S6 terminal parameters put
\begin{equation}
 \lambda(\ell_0,\ell_1,\ell_2)
 =12\ell_0-4\ell_1-3\ell_2                 \tag{1.3}
\end{equation}
and define the abstract terminal group
\begin{equation}
 G(\ell_0,\ell_1,\ell_2)=
 \left\langle c,x,y\ \middle|\
 c\text{ central},\ xy=c^{\ell_0},\
 x^3=c^{\ell_1},\ y^4=c^{\ell_2}
 \right\rangle .                            \tag{1.4}
\end{equation}
After eliminating \(y\), and using the ordered generator basis \((x,c)\),
its ordered row-relation matrix is
\begin{equation}
 \mathsf R(\ell_0,\ell_1,\ell_2)=
 \begin{pmatrix}
 3&-\ell_1\\
 4&\ell_2-4\ell_0
 \end{pmatrix},
 \qquad
 \det\mathsf R=-\lambda(\ell_0,\ell_1,\ell_2).              \tag{1.5}
\end{equation}
The letter \(p\) in this note always denotes the Erd\H{o}s--Straus prime.
The signed terminal value is denoted
\[
 q_\eps=p^\eps\in\{1,p\},\qquad \eps\in\{0,1\}.
\]
Thus no identification of the external prime \(p\) with the terminal value
is made in the \(\eps=0\) channel.

For any integral \(2\times2\) matrix \(A\), define its row-relation lattice
and row-presented group by
\[
 \Rel(A)=\{v^{\mathsf T}A:v\in\ZZ^2\}\subseteq\ZZ^2,
 \qquad \mathcal G(A)=\ZZ^2/\Rel(A)=\coker(A^{\mathsf T}).    \tag{1.6}
\]
This convention makes every matrix direction in the bridge explicit.

\subsection{The complete first-half shell selector}

\begin{lemma}[first-half denominator]\label{lem:esksb-first-half}
Let \(p\equiv1\pmod4\) be an odd prime.  If
\((x,y,z)\in\mathscr S(p)\) and \(x\leq y\leq z\), then
\[
 \frac p4<x\leq\frac p2.
\]
\end{lemma}

\begin{proof}
Positivity of the two reciprocals other than \(1/x\) gives
\(1/x<4/p\), hence \(x>p/4\).  Suppose \(x>p/2\).  Clearing denominators
gives
\[
 4xyz=p(xy+xz+yz),
\]
so \(p\mid xyz\).  If \(p\mid x\), then \(x,y,z\geq p\) and the reciprocal
sum is at most \(3/p\).  If \(p\mid y\), then \(y,z\geq p\), while
\(1/x<2/p\), so the sum is strictly less than \(4/p\).  Both cases are
impossible.  Therefore \(z=kp\) for some \(k\in\ZZpos\), while neither
\(x\) nor \(y\) is divisible by \(p\).  Substitution yields
\[
 (4k-1)xy=kp(x+y).
\]
Reduction modulo \(p\) gives \(p\mid4k-1\).  Write \(4k-1=rp\).
Since \(p\equiv1\pmod4\), one has \(r\equiv3\pmod4\), hence \(r\geq3\).
After cancelling \(p\),
\[
 rxy=k(x+y).
\]
As \(x\leq y\), division by \(xy\) gives
\[
 x\leq\frac{2k}{r}
 =\frac{rp+1}{2r}
 =\frac p2+\frac1{2r}<\frac{p+1}{2}.
\]
The integer \(x>p/2\) must be at least \((p+1)/2\), a contradiction.
\end{proof}

For fixed \(j\in J_p\), define the ordered residual shell
\begin{equation}
 \mathscr P_{p,j}=
 \{(y,z)\in\ZZpos^2:R_jyz=N_j(y+z)\}                         \tag{2.1}
\end{equation}
and the oriented divisor shell
\begin{equation}
 \mathscr D_{p,j}=
 \{d\in\ZZpos:d\mid N_j^2,\ d\equiv-N_j\pmod {R_j}\}.       \tag{2.2}
\end{equation}

\begin{proposition}[oriented shell bijection]\label{prop:esksb-shell-bijection}
The maps
\begin{align*}
 \delta_{p,j}:\mathscr P_{p,j}&\longrightarrow\mathscr D_{p,j},
 &(y,z)&\longmapsto R_jy-N_j,\\
 \rho_{p,j}:\mathscr D_{p,j}&\longrightarrow\mathscr P_{p,j},
 &d&\longmapsto
 \left(\frac{N_j+d}{R_j},
       \frac{N_j+N_j^2/d}{R_j}\right)
\end{align*}
are mutually inverse.  The involution
\[
 \iota_{N_j}(d)=\frac{N_j^2}{d}
\]
corresponds exactly to \((y,z)\mapsto(z,y)\).
\end{proposition}

\begin{proof}
Abbreviate \(R=R_j\) and \(N=N_j\).  Equation~(2.1) is equivalent to
\begin{equation}
 (Ry-N)(Rz-N)=N^2.                                            \tag{2.3}
\end{equation}
Both factors are positive: for example, division of
\(Ryz=N(y+z)\) by \(z\) gives \(Ry=N(1+y/z)>N\), and similarly for
\(Rz-N\).  Therefore \(d=Ry-N\) is a positive divisor of \(N^2\) and
\(d\equiv-N\pmod R\).

Conversely, let \(d\in\mathscr D_{p,j}\) and put \(e=N^2/d\).  By~(1.2),
\(d\) and \(N\) are units modulo \(R\); hence
\[
 e\equiv N^2(-N)^{-1}\equiv-N\pmod R.
\]
Thus both components of \(\rho_{p,j}(d)\) are positive integers.  Their
 two shifted factors are \(d\) and \(e\), so~(2.3), and therefore~(2.1),
 holds.  The formulas visibly compose to the identity.  Replacing \(d\) by
\(e=N^2/d\) exchanges the two reconstructed components.
\end{proof}

For \(\eps\in\{0,1\}\), define the fixed-shell selector and the fully marked
selector by
\begin{align}
 \mathcal W_\eps(p,j)
 &=\{u\in\ZZpos:u\mid a_j^2,\ R_j\mid4u+p^\eps\},             \tag{2.4}\\
 \mathcal W_\eps(p)
 &=\{(j,u):j\in J_p,\ u\in\mathcal W_\eps(p,j)\}.           \tag{2.5}
\end{align}
The tag \(\eps=0\) is the exterior, or \(E\), channel; \(\eps=1\) is the
middle, or \(M\), channel.
Whenever a disjoint union of the two channels is written below, it means the
tagged coproduct
\[
 \{0\}\times\mathcal W_0(p)\ \sqcup\
 \{1\}\times\mathcal W_1(p).
\]
This convention is necessary because the two raw subsets can intersect.

\begin{proposition}[exact two-channel decomposition]\label{prop:esksb-two-channel}
The maps
\begin{align*}
 \iota_E:\mathcal W_0(p,j)&\longrightarrow\mathscr D_{p,j},
 &u&\longmapsto p^2u,\\
 \iota_M:\mathcal W_1(p,j)&\longrightarrow\mathscr D_{p,j},
 &u&\longmapsto pu
\end{align*}
are bijections onto the subsets on which the \(p\)-adic valuations are \(2\)
and \(1\), respectively.  The valuation-zero subset is
\(\iota_{N_j}(\iota_E(\mathcal W_0(p,j)))\).  Consequently
\begin{equation}
 \mathscr D_{p,j}=
 \iota_E(\mathcal W_0(p,j))\ \sqcup\
 \iota_M(\mathcal W_1(p,j))\ \sqcup\
 \iota_{N_j}(\iota_E(\mathcal W_0(p,j))),                    \tag{2.6}
\end{equation}
and
\begin{equation}
 |\mathscr D_{p,j}|=2|\mathcal W_0(p,j)|+|\mathcal W_1(p,j)|. \tag{2.7}
\end{equation}
\end{proposition}

\begin{proof}
Because \(p\nmid a_j\), every divisor \(d\mid p^2a_j^2\) has
\(v_p(d)\in\{0,1,2\}\).

If \(d=p^2u\), then \(u\mid a_j^2\).  The shell congruence is
\[
 R_j\mid d+N_j=p(pu+a_j).
\]
After cancelling the unit \(p\) modulo \(R_j\) and multiplying by \(4\),
the identity \(4a_j=p+R_j\) makes this equivalent to
\(R_j\mid4u+1\).  This proves the \(E\)-channel assertion in both
directions.  If \(d=pu\), the same calculation gives
\[
 R_j\mid d+N_j
 \Longleftrightarrow R_j\mid u+a_j
 \Longleftrightarrow R_j\mid4u+p,
\]
which proves the \(M\)-channel assertion.  Complementation preserves the
shell congruence by Proposition~\ref{prop:esksb-shell-bijection}, exchanges
valuations \(0\) and \(2\), and preserves valuation \(1\).  This proves the
disjoint decomposition and its cardinality formula.
\end{proof}

For \(u\in\mathcal W_\eps(p,j)\), define the two exact quotients
\begin{equation}
 T=\frac{4u+p^\eps}{R_j},
 \qquad K=\frac{p^{1-\eps}T+1}{4}.                            \tag{2.8}
\end{equation}
Since \(R_j\equiv3\pmod4\) and \(p^\eps\equiv1\pmod4\), the first equation
gives \(T\equiv3\pmod4\).  Hence \(K\in\ZZpos\).  Direct calculation,
without dropping the tag, gives
\begin{equation}
 R_jK=a_j+p^{1-\eps}u,
 \qquad Ta_j=u+p^\eps K.                                    \tag{2.9}
\end{equation}
Indeed, multiplication by \(4\) changes these identities into
\begin{align*}
 4R_jK
 &=p^{1-\eps}(4u+p^\eps)+R_j
 =4p^{1-\eps}u+p+R_j,\\
 4Ta_j
 &=T(p+R_j)
 =pT+4u+p^\eps
 =4u+4p^\eps K.
\end{align*}

\begin{proposition}[denominator reconstruction]\label{prop:esksb-reconstruct}
The map
\begin{equation}
 \Psi_{p,j,\eps}:\mathcal W_\eps(p,j)\longrightarrow\mathscr S(p),
 \qquad
 u\longmapsto
 \left(a_j,\ pK,\ \frac{p^\eps a_jK}{u}\right)             \tag{2.10}
\end{equation}
is well-defined.  Its two oriented shell factors are
\begin{equation}
 R_j(pK)-N_j=p^{2-\eps}u,
 \qquad
 R_j\left(\frac{p^\eps a_jK}{u}\right)-N_j
 =p^\eps\frac{a_j^2}{u}.                                    \tag{2.11}
\end{equation}
Moreover,
\begin{equation}
 \ES(p)\quad\Longleftrightarrow\quad
 \bigl(\{0\}\times\mathcal W_0(p)\bigr)\sqcup
 \bigl(\{1\}\times\mathcal W_1(p)\bigr)\neq\varnothing.   \tag{2.12}
\end{equation}
\end{proposition}

\begin{proof}
The last coordinate in~(2.10) is integral.  To see this without assuming it,
reduce the second identity in~(2.9) modulo \(u\), multiply by \(a_j\), and
use \(u\mid a_j^2\): this gives \(u\mid p^\eps a_jK\).
The two identities~(2.11) now follow from the first identity in~(2.9).
Their product is \(N_j^2\), so Proposition~\ref{prop:esksb-shell-bijection}
shows that~(2.10) lies in \(\mathscr S(p)\).

Conversely, suppose \(\mathscr S(p)\neq\varnothing\), and sort a solution.
Lemma~\ref{lem:esksb-first-half} shows that its smallest denominator is uniquely
of the form \(x=a_j\) with \(j\in J_p\): explicitly
\[
 R_j=4x-p\equiv3\pmod4,
 \qquad j=\frac{4x-p-3}{4}.
\]
Proposition~\ref{prop:esksb-shell-bijection} supplies a divisor
\(d\in\mathscr D_{p,j}\), and Proposition~\ref{prop:esksb-two-channel} places
either \(d\) or its complement in the image of \(\iota_E\), or places \(d\)
in the image of \(\iota_M\).  Hence one of the two marked selector sets is
nonempty.  The forward implication follows from~(2.10).
\end{proof}

\subsection{The exact selector--terminal-lattice equivalence}

Define the normalized, positively oriented, divisor-supported terminal slice
\begin{equation}
 \mathcal H_\eps(p)=
 \left\{(j,\ell_0,\ell_1)\in\ZZ^3:
 \begin{array}{l}
 j\in J_p,\quad \lambda(\ell_0,\ell_1,1)=p^\eps,\\
 T(\ell_0):=4\ell_0-1>0,\\
 u(j,\ell_0,\ell_1):=\ell_1+jT(\ell_0)>0,\\
 u(j,\ell_0,\ell_1)\mid a_j^2
 \end{array}
 \right\}.                                                   \tag{3.1}
\end{equation}
Let \(\mathcal H_\eps(p,j)\) be the fibre with the displayed marker \(j\).

For a selector state define its relation matrix
\begin{equation}
 \mathsf A_\eps(p;j,u)=
 \begin{pmatrix}
 R_j&-u\\
 4&-T
 \end{pmatrix},
 \qquad R_jT=4u+p^\eps.                                     \tag{3.2}
\end{equation}
For the marker \(j\), retain both inverse unimodular shears
\begin{equation}
 \mathsf U_j=
 \begin{pmatrix}1&-j\\0&1\end{pmatrix},
 \qquad
 \mathsf L_j=\mathsf U_j^{-1}=
 \begin{pmatrix}1&j\\0&1\end{pmatrix}.                    \tag{3.3}
\end{equation}

\begin{theorem}[strong marked bridge theorem]\label{thm:esksb-bridge}
For every \(\eps\in\{0,1\}\), the map
\begin{equation}
 \Theta_{\eps,p}:\mathcal W_\eps(p)\longrightarrow\mathcal H_\eps(p),
 \qquad
 (j,u)\longmapsto
 \left(j,\frac{T+1}{4},u-jT\right)                          \tag{3.4}
\end{equation}
is a bijection.  Its inverse is
\begin{equation}
 \Theta_{\eps,p}^{-1}(j,\ell_0,\ell_1)
 =\left(j,\ell_1+j(4\ell_0-1)\right).                       \tag{3.5}
\end{equation}
For corresponding points, with
\(\ell=(\ell_0,\ell_1,1)\), one has the exact matrix identities
\begin{equation}
 \boxed{
 \mathsf U_j\mathsf A_\eps(p;j,u)
 =\mathsf R(\ell_0,\ell_1,1),
 \qquad
 \mathsf A_\eps(p;j,u)
 =\mathsf L_j\mathsf R(\ell_0,\ell_1,1).}                  \tag{3.6}
\end{equation}
In particular,
\begin{equation}
 \Rel(\mathsf A_\eps(p;j,u))
 =\Rel(\mathsf R(\ell_0,\ell_1,1)),
 \qquad
 \det\mathsf A_\eps=\det\mathsf R=-p^\eps.                \tag{3.7}
\end{equation}
Thus the bridge intertwines the entire integral row presentation, not only
its determinant.
\end{theorem}

\begin{proof}
Take \((j,u)\in\mathcal W_\eps(p)\).  Section~2 proves that \(T\) is a
positive integer satisfying \(T\equiv3\pmod4\), so
\(\ell_0=(T+1)/4\in\ZZ\).  Put \(\ell_1=u-jT\).  Since
\(R_j=4j+3\),
\begin{align*}
 \lambda(\ell_0,\ell_1,1)
 &=3(T+1)-4(u-jT)-3\\
 &=(4j+3)T-4u=p^\eps.
\end{align*}
The positivity and divisor-support conditions in~(3.1) hold because
\(u=\ell_1+jT\).  Thus~(3.4) is well-defined.

Conversely, take \((j,\ell_0,\ell_1)\in\mathcal H_\eps(p)\), and set
\[
 T=4\ell_0-1,
 \qquad u=\ell_1+jT.
\]
Expanding \(\lambda(\ell_0,\ell_1,1)=p^\eps\) gives
\[
 (4j+3)T-4u=p^\eps,
\]
so \(R_j\mid4u+p^\eps\).  All other selector conditions occur explicitly
in~(3.1).  The two constructions recover \(T,u,\ell_0,\ell_1\) literally,
which proves bijectivity and~(3.5).

Finally, direct multiplication gives
\[
 \begin{pmatrix}1&-j\\0&1\end{pmatrix}
 \begin{pmatrix}4j+3&-u\\4&-T\end{pmatrix}
 =
 \begin{pmatrix}3&-u+jT\\4&-T\end{pmatrix}
 =
 \begin{pmatrix}3&-\ell_1\\4&1-4\ell_0\end{pmatrix}.
\]
Both shears have determinant \(1\), so their integral row operations are
invertible and the row lattices are equal.  The signed determinant is
\(-R_jT+4u=-p^\eps\), proving~(3.7).
\end{proof}

\begin{corollary}[presented groups and the Erd\H{o}s--Straus equation]
For every corresponding selector/S6 point,
\begin{equation}
 \mathcal G(\mathsf A_\eps(p;j,u))
 =\mathcal G(\mathsf R(\ell_0,\ell_1,1))
 \cong G(\ell_0,\ell_1,1)
 \cong\ZZ/p^\eps\ZZ.                                       \tag{3.8}
\end{equation}
Here \(\ZZ/1\ZZ\) is the trivial group.  Moreover the direct map
\begin{equation}
 \widehat\Psi_{\eps,p}:\mathcal H_\eps(p)\longrightarrow\mathscr S(p)
                                                                    \tag{3.9}
\end{equation}
is obtained by~(3.5), (2.8), and
\begin{equation}
 (j,\ell_0,\ell_1)\longmapsto
 \left(a_j,\ pK,\ \frac{p^\eps a_jK}{u}\right),
 \quad
 \begin{cases}
 T=4\ell_0-1,\\
 u=\ell_1+jT,\\
 K=(p^{1-\eps}T+1)/4.
 \end{cases}                                                \tag{3.10}
\end{equation}
It satisfies
\begin{equation}
 \boxed{\quad
 \ES(p)\Longleftrightarrow
 \bigl(\{0\}\times\mathcal H_0(p)\bigr)\sqcup
 \bigl(\{1\}\times\mathcal H_1(p)\bigr)\neq\varnothing.
 \quad}                                                       \tag{3.11}
\end{equation}
\end{corollary}

\begin{proof}
The relation \(xy=c^{\ell_0}\) in~(1.4) gives
\(y=x^{-1}c^{\ell_0}\); since \(c\) is central, \(x\) and \(y\) commute.
The remaining two relations are the rows of~(1.5).  The greatest common
divisor of the four entries of \(\mathsf R\) is \(1\), because it divides
both \(3\) and \(4\).  Its Smith form is therefore
\(\operatorname{diag}(1,p^\eps)\), proving~(3.8).  Formula~(3.10) is the
composition of Theorem~\ref{thm:esksb-bridge} with
Proposition~\ref{prop:esksb-reconstruct}, and~(3.11) follows from~(2.12).
\end{proof}

\begin{remark}[the selector tag and the solution map]
The channel tag is not redundant.  For \(p=13\), \(j=0\), and \(u=2\),
one has \(u\in\mathcal W_0(p,j)\cap\mathcal W_1(p,j)\), but the two tags
give different terminal values and different marked terminal points.  Nor is
the map to unordered denominator triples injective: in the \(M\)-channel at
\(p=13,j=0\), the complementary divisors \(u=2\) and \(u=8=a_0^2/2\)
both give the unordered solution \(\{4,26,52\}\).  The bijection in
Theorem~\ref{thm:esksb-bridge} is a bijection of marked selector states with marked
terminal states, not a bijection with unordered solutions.
\end{remark}

\begin{corollary}[fixed-fibre affine coordinates and the intrinsic gate]
Put
\[
 C_\eps=\frac{p^\eps+3}{4},
 \qquad D_{\eps,j}=j+C_\eps=\frac{R_j+p^\eps}{4}.
\]
Every integral point with \(\ell_2=1\) and
\(\lambda(\ell_0,\ell_1,1)=p^\eps\) is uniquely
\begin{equation}
 (\ell_0,\ell_1,\ell_2)=(n,3n-C_\eps,1),\qquad n\in\ZZ,      \tag{3.12}
\end{equation}
and, for each fixed marker \(j\in J_p\), define
\begin{equation}
 T=4n-1,
 \qquad u=R_jn-D_{\eps,j},
 \qquad K=\frac{p^{1-\eps}T+1}{4}.                          \tag{3.13}
\end{equation}
Then \(T\equiv3\pmod4\), so \(K\in\ZZ\), and the congruence equation
\(R_jT=4u+p^\eps\) holds identically.  On the locus \(T>0\) and
\(u>0\), one has
\begin{equation}
 \gcd(R_j,u)=1,
 \qquad
 u\mid a_j^2\Longleftrightarrow u\mid K^2.                  \tag{3.14}
\end{equation}
Thus the nonlinear divisor gate can be written entirely in the S6
coordinates, but it does not disappear.
\end{corollary}

\begin{proof}
Equation~(3.12) is the direct solution of
\(12\ell_0-4\ell_1-3=p^\eps\).  Substitution gives the first two
formulas in~(3.13), and \(T\equiv3\pmod4\) proves that the displayed
integer \(K\) is well defined.  The definitions give
\[
 4u=4R_jn-(R_j+p^\eps)=R_jT-p^\eps,
\]
so the congruence equation holds for every integral point in~(3.12).
If a positive integer divides both \(R_j\) and \(u\), then the congruence
equation makes it divide \(p^\eps\); (1.2) therefore gives
\(\gcd(R_j,u)=1\).  Without invoking the selector hypothesis, direct
substitution gives
\[
 4\bigl(R_jK-p^{1-\eps}u-a_j\bigr)
 =p^{1-\eps}(R_jT-4u)-p=0.
\]
Thus \(a_j\equiv R_jK\pmod u\).  Since \(R_j\) is a unit modulo \(u\),
\(a_j^2\equiv R_j^2K^2\pmod u\), which proves (3.14) in both directions.
\end{proof}

The normalization automatically lies in the admissible residue class used by
the S6 terminal presentation:
\begin{equation}
 \ell_2=1\text{ is odd},
 \qquad \ell_1\equiv2\pmod3.                                \tag{3.15}
\end{equation}
Indeed reduction of \(\lambda(\ell_0,\ell_1,1)=p^\eps\equiv1\pmod3\)
gives \(-\ell_1\equiv1\pmod3\).

\begin{remark}[orientation and the constructed S6 candidate]
The geometrically reconstructed candidate elsewhere in the audited corpus
uses the oriented terminal parameter
\(\lambda(0,1,-1)=-1\).  The two arithmetic channels in this note use
the positive terminal values \(+1\) and \(+p\).  No audited map identifies
the \(-1\) terminal datum with either positive datum.  Consequently the row
shear proved here does not, by itself, place any selector state on the
particular complex threefold constructed in the \(-1\) branch.
\end{remark}

\subsection{The exact remaining arithmetic incidence}

The bridge can be expressed without an untyped phrase such as ``some divisor
of a square.''  The next factor coordinates are canonical and retain their
orientation.

\begin{lemma}[canonical ordered factor coordinates]\label{lem:esksb-factors}
For \(a,u\in\ZZpos\), the condition \(u\mid a^2\) is equivalent to the
existence of a unique ordered triple \((r,s,t)\in\ZZpos^3\) such that
\begin{equation}
 u=r^2s,
 \qquad a=rst,
 \qquad \gcd(r,t)=1.                                        \tag{4.1}
\end{equation}
Explicitly, with \(g=\gcd(a,u)\),
\begin{equation}
 r=\frac ug,
 \qquad s=\frac{g^2}{u},
 \qquad t=\frac ag.                                         \tag{4.2}
\end{equation}
\end{lemma}

\begin{proof}
Write \(v_q(a)=\alpha\) and \(v_q(u)=\beta\leq2\alpha\) prime by prime.
Then
\[
 v_q(g)=\min(\alpha,\beta),
 \qquad
 v_q(g^2/u)=2\min(\alpha,\beta)-\beta\geq0.
\]
Thus~(4.2) consists of positive integers.  If \(\beta\leq\alpha\), then
\(q\nmid r\); if \(\beta\geq\alpha\), then \(q\nmid t\).  Hence
\(\gcd(r,t)=1\), and direct multiplication gives~(4.1).

Conversely,~(4.1) gives
\(a^2/u=s t^2\in\ZZpos\).  It also gives
\(\gcd(a,u)=rs\), because \(\gcd(r,t)=1\).  Formula (4.2) then recovers
\((r,s,t)\), proving uniqueness.
\end{proof}

\begin{proposition}[factor form of both selector channels]\label{prop:esksb-factor-congruence}
For the unique factorization
\(u=r^2s\), \(a_j=rst\), \(\gcd(r,t)=1\), one has
\begin{align}
 u\in\mathcal W_0(p,j)
 &\Longleftrightarrow R_j\mid pr+t,                           \tag{4.3}\\
 u\in\mathcal W_1(p,j)
 &\Longleftrightarrow R_j\mid r+t.                            \tag{4.4}
\end{align}
\end{proposition}

\begin{proof}
All of \(r,s,t\) are units modulo \(R_j\), because their product is \(a_j\)
and \(\gcd(a_j,R_j)=1\).  For the \(E\)-channel,
\[
 t(4u+1)=4r^2st+t=(p+R_j)r+t=(pr+t)+R_jr.
\]
Cancelling \(t\) modulo \(R_j\) proves~(4.3).  For the \(M\)-channel,
\[
 4u+p\equiv4(u+a_j)=4rs(r+t)\pmod {R_j},
\]
and cancellation of the unit \(4rs\) proves~(4.4).
\end{proof}

\begin{theorem}[necessary-and-sufficient counterexample criterion]
A prime \(p\equiv1\pmod {12}\) is a counterexample to the
Erd\H{o}s--Straus equation if and only if, for every \(j\in J_p\) and every
ordered \((r,s,t)\in\ZZpos^3\) satisfying
\[
 a_j=rst,
 \qquad \gcd(r,t)=1,
\]
one has simultaneously
\begin{equation}
 R_j\nmid pr+t,
 \qquad R_j\nmid r+t.                                       \tag{4.5}
\end{equation}
Equivalently, with the tag retained, it is necessary and sufficient that
\begin{equation}
 R_j\nmid p^{1-\eps}r+t                                     \tag{4.6}
\end{equation}
for every \(j\in J_p\), every \(\eps\in\{0,1\}\), and every such factor
triple.
\end{theorem}

\begin{proof}
Lemma~\ref{lem:esksb-factors} is a bijection between all allowed divisors in
the selector~(2.4) and all displayed ordered factor triples.  Proposition
\ref{prop:esksb-factor-congruence} translates the two selector congruences without
loss.  The assertion now follows directly from~(2.12), or equivalently
(3.11), by negating the finite existential statement.
\end{proof}

This theorem isolates the unresolved problem exactly: the S6 row shear has
not removed the bounded factor-incidence condition~(4.5).

\subsection{Why the labels and support cannot be forgotten}

\begin{lemma}[full terminal fibres are affine lattices]
The kernel of \(\lambda:\ZZ^3\to\ZZ\) is
\begin{equation}
 \ker\lambda=\ZZ(1,3,0)\oplus\ZZ(1,0,4).                    \tag{5.1}
\end{equation}
If \(q\equiv1\pmod {12}\) and \(m_q=(q+11)/12\), then
\begin{equation}
 \lambda^{-1}(q)=
 \{(m_q+r+s,\ 2+3r,\ 1+4s):r,s\in\ZZ\}.                   \tag{5.2}
\end{equation}
In the normalized slice \(\ell_2=1\), this is an infinite affine line.
\end{lemma}

\begin{proof}
If \(12a-4b-3c=0\), reduction modulo \(4\) gives \(c=4s\).  The remaining
equation gives \(b=3(a-s)\); putting \(r=a-s\) yields
\((a,b,c)=r(1,3,0)+s(1,0,4)\), uniquely.  Since
\(\lambda(m_q,2,1)=12m_q-11=q\), translation by~(5.1) proves~(5.2).
\end{proof}

The lemma proves that the mere existence of a terminal parameter with signed
value \(1\) or \(p\) is automatic.  It cannot imply selector nonemptiness.
The following examples show separately why the marker and the divisor-support
condition are both part of the theorem's domain.

\begin{example}[forgetting \(j\) and \(p\) is noninjective]\label{ex:esksb-tags}
For each eligible prime \(q\), define the fixed-prime forgetful map
\[
 \mathfrak f_q:\mathcal H_0(q)\longrightarrow\ZZ^3,
 \qquad
 \mathfrak f_q(j,\ell_0,\ell_1)=(\ell_0,\ell_1,1).
\]
At \(p=13\), the two \(E\)-states
\[
 (j,u,T)=(0,2,3),\qquad(1,5,3)
\]
satisfy
\[
 \Theta_{0,13}(0,2)=(0,1,2),
 \qquad
 \Theta_{0,13}(1,5)=(1,1,2),
\]
and both images are sent by \(\mathfrak f_{13}\) to
\((1,2,1)\).  To compare different prime fibres, form the tagged coproduct
and the global forgetful map
\[
 \mathfrak F:
 \bigsqcup_{\substack{q\ \text{prime}\\q\equiv1\pmod {12}}}
 \bigl(\{q\}\times\mathcal H_0(q)\bigr)
 \longrightarrow\ZZ^3,
 \qquad
 \mathfrak F(q;j,\ell_0,\ell_1)=(\ell_0,\ell_1,1).
\]
The state \((p,j,u)=(73,1,5)\) has
\(\Theta_{0,73}(1,5)=(1,1,2)\), so its tagged image also maps under
\(\mathfrak F\) to \((1,2,1)\).  Therefore neither \(j\) nor the external
prime \(p\) is recoverable after applying the explicitly displayed
forgetful maps.
\end{example}

\begin{example}[ambient positivity does not imply divisor support]\label{ex:esksb-ambient}
For \(p=13\), \(j=2\), one has \(R_j=11\) and \(a_j=6\).  The positive
normalized \(E\)-line has
\[
 (\ell_0,\ell_1)=(n,3n-1),
 \qquad u=11n-3,
\]
and the positive normalized \(M\)-line has
\[
 (\ell_0,\ell_1)=(n,3n-4),
 \qquad u=11n-6.
\]
Both lines are infinite.  The only positive \(u\leq a_j^2=36\) are
\[
 E:\ 8,19,30,
 \qquad M:\ 5,16,27,
\]
and none divides \(36\).  Hence
\(\mathcal H_0(13,2)=\mathcal H_1(13,2)=\varnothing\), even though both
ambient positive S6 lines are nonempty.  The prime \(13\) is solved at other
markers, so this is a fixed-shell counterexample to forgetting support, not
an Erd\H{o}s--Straus counterexample.
\end{example}

\subsection{The natural unmarked \texorpdfstring{\(S^6\)}{S6} point and its complete divisor gate}

The most direct attempted use of the S6 presentation starts from one point in
the signed \(p\)-fibre and shears it through every shell.  Its exact content
can be classified completely.

Write
\begin{equation}
 p=12m-11,
 \qquad m=h+1\geq2,
 \qquad J_p=\{0,\ldots,3m-4\}.                              \tag{6.1}
\end{equation}
Retain the unmarked S6 point
\begin{equation}
 \ell^{\mathrm{aff}}=(m,2,1),
 \qquad \lambda(\ell^{\mathrm{aff}})=p,                     \tag{6.2}
\end{equation}
and for \(j\in J_p\) put
\begin{equation}
 T_0=4m-1,
 \quad u_j=2+jT_0,
 \quad a_j=3m+j-2,
 \quad R_j=4j+3.                                             \tag{6.3}
\end{equation}
Then
\begin{equation}
 R_jT_0=4u_j+p,
 \qquad R_jm=a_j+u_j,
 \qquad T_0a_j-u_j=pm,                                      \tag{6.4}
\end{equation}
and
\begin{equation}
 \mathsf U_j
 \begin{pmatrix}R_j&-u_j\\4&-T_0\end{pmatrix}
 =
 \begin{pmatrix}3&-2\\4&1-4m\end{pmatrix}
 =\mathsf R(m,2,1).                                         \tag{6.5}
\end{equation}

\begin{theorem}[complete affine-line gate]\label{thm:esksb-affine-gate}
For prime \(p=12m-11\), \(m\geq2\), and \(0\leq j<3m-3\),
\begin{equation}
 u_j\mid a_j^2
 \Longleftrightarrow
 \bigl(j=0\text{ and }2\mid m\bigr)
 \ \text{or}\ 
 \bigl(j=2\text{ and }8\mid m\bigr).                       \tag{6.6}
\end{equation}
Consequently the single unmarked point~(6.2), with all markers restored,
produces an Erd\H{o}s--Straus witness if and only if
\(p\equiv13\pmod {24}\).  In that case \(j=0,u=2,K=m\) gives
\begin{equation}
 \frac4p=
 \frac1{3m-2}
 +\frac1{pm}
 +\frac{1}{pm(3m-2)/2}.                                    \tag{6.7}
\end{equation}
For \(p\equiv1\pmod {24}\), this affine S6 point supplies no supported
selector state at any \(j\in J_p\).
\end{theorem}

\begin{proof}
The identities~(6.4) and~(6.5) follow by direct expansion.  Assume first
that \(u_j\mid a_j^2\), and put \(g=\gcd(u_j,a_j)\).  Write
\(u_j=gu'\) and \(a_j=ga'\) with \(\gcd(u',a')=1\).  The assumed
divisibility gives \(u'\mid g(a')^2\), hence \(u'\mid g\), and therefore
\begin{equation}
 u_j\mid g^2.                                                \tag{6.8}
\end{equation}
The last identity in~(6.4) gives \(g\mid pm\).  Since \(0<a_j<p\), the
primality of \(p\) and \(g\mid a_j\) imply \(p\nmid g\); therefore
\(g\mid m\).  Reducing \(u_j=2+j(4m-1)\) modulo \(m\) now gives
\begin{equation}
 g\mid j-2.                                                  \tag{6.9}
\end{equation}

Suppose \(j\neq0,2\).  Then \(j\geq1\), and~(6.8)--(6.9) imply
\(u_j\leq g^2\leq(j-2)^2\).  But
\begin{equation}
 u_j-(j-2)^2=j(4m+3-j)-2>0                                  \tag{6.10}
\end{equation}
For \(m\geq2\) and \(1\leq j<3m-3\), indeed
\(4m+3-j\geq m+7\), so the right-hand side of~(6.10) is at least
\(m+5>0\).  This is a contradiction.  Thus only \(j=0\) or \(j=2\)
can occur.

At \(j=0\), \(u_0=2\) and \(a_0=3m-2\), so
\(u_0\mid a_0^2\) exactly when \(m\) is even.  At \(j=2\),
\(u_2=8m\) and \(a_2=3m\), so
\[
 u_2\mid a_2^2
 \Longleftrightarrow 8m\mid9m^2
 \Longleftrightarrow 8\mid m.
\]
These calculations also prove both converse implications in~(6.6).

Now \(p=12m-11\) is congruent to \(13\pmod {24}\) exactly when \(m\) is
even.  The \(j=0\) state lies in the \(M\)-channel, and~(2.8) gives
\(T=T_0\) and \(K=(T_0+1)/4=m\).  Formula~(2.10) is precisely~(6.7).
If \(m\) is odd,~(6.6) excludes every marker on this particular S6 point.
\end{proof}

The last sentence is deliberately local to~(6.2).  For example, \(p=73\)
has \(m=7\) odd, so the affine point contributes nothing; nevertheless
\((j,u,T;\ell_0,\ell_1)=(1,1,11;3,-10)\) is an \(M\)-state.  Thus failure of
the affine point is not failure of the full selector.

\subsection{Exact descendant maps transported through the bridge}

The bridge also transports a verified selector witness along multiplicative
descendant rays.  This is a map from an already nonempty domain, not a
nonemptiness theorem.

Fix a supported state at \((p,j)\), retain \(R=R_j\), and let \(b\in\ZZpos\).
Define
\begin{equation}
 a_b=ba_j,
 \qquad p_b=4a_b-R=bp+(b-1)R.                               \tag{7.1}
\end{equation}
Whenever \(p_b\) is again a prime congruent to \(1\pmod {12}\), the same
marker \(j\) belongs to \(J_{p_b}\), since \(R<p\leq p_b\), and the following
maps have the displayed domains and codomains:
\begin{align}
 \Phi_{b;p}^E:\mathcal H_0(p,j)&\longrightarrow\mathcal H_0(p_b,j),
 &(j,\ell_0,\ell_1)&\longmapsto(j,\ell_0,\ell_1),             \tag{7.2}\\
 \Phi_{b;p}^M:\mathcal H_1(p,j)&\longrightarrow\mathcal H_1(p_b,j),
 &(j,\ell_0,\ell_1)&\longmapsto
 (j,b\ell_0,b\ell_1-j(b-1)).                               \tag{7.3}
\end{align}
Write
\[
 \theta_{\eps,p,j}:\mathcal W_\eps(p,j)\longrightarrow
 \mathcal H_\eps(p,j),
 \qquad
 \theta_{\eps,p,j}(u)=\Theta_{\eps,p}(j,u)
\]
for the fixed-marker restriction of the marked bridge.

\begin{proposition}[intertwining of descendant actions]
The selector maps corresponding to~(7.2) and~(7.3) are the typed maps
\begin{align*}
 \psi_{b;p}^E:\mathcal W_0(p,j)&\longrightarrow\mathcal W_0(p_b,j),
 &u&\longmapsto u,\\
 \psi_{b;p}^M:\mathcal W_1(p,j)&\longrightarrow\mathcal W_1(p_b,j),
 &u&\longmapsto bu.
\end{align*}
They intertwine the fixed-marker bridges exactly:
\begin{equation}
 \Phi_{b;p}^\bullet\circ\theta_{\eps,p,j}
 =\theta_{\eps,p_b,j}\circ\psi_{b;p}^\bullet,
 \qquad
 (\bullet,\eps)\in\{(E,0),(M,1)\}.                         \tag{7.4}
\end{equation}
Equivalently,
\[
 \theta_{0,p_b,j}^{-1}\Phi_{b;p}^E\theta_{0,p,j}(u)=u,
 \qquad
 \theta_{1,p_b,j}^{-1}\Phi_{b;p}^M\theta_{1,p,j}(u)=bu.
\]
On the accompanying arithmetic metadata these maps replace
\((a_j,p,u)\) by \((a_b,p_b,u)\) and \((a_b,p_b,bu)\), respectively.
For the latter,
\begin{equation}
 T_b=bT+b-1,
 \qquad u_b=bu,
 \qquad K_b=bK,                                              \tag{7.5}
\end{equation}
and
\begin{equation}
 \lambda(b\ell_0,b\ell_1-j(b-1),1)=p_b.                    \tag{7.6}
\end{equation}
For either tag the maps are injective.  For \(c\in\ZZpos\), put
\[
 a_{bc}=bc\,a_j,
 \qquad
 p_{bc}=4a_{bc}-R=bc\,p+(bc-1)R.
\]
Whenever \(p_b\) and \(p_{bc}\) are primes congruent to \(1\pmod {12}\),
the displayed maps are composable and satisfy
\begin{equation}
 \Phi_{c;p_b}^\bullet\circ\Phi_{b;p}^\bullet
 =\Phi_{bc;p}^\bullet                                      \tag{7.7}
\end{equation}
\end{proposition}

\begin{proof}
In the \(E\)-channel, \(u\mid a_j^2\) implies \(u\mid a_b^2\), and the
congruence \(R\mid4u+1\) and quotient \(T\) are unchanged.  Hence
\(\ell_0,\ell_1\) are unchanged by~(3.4).

In the \(M\)-channel, \(bu\mid b^2a_j^2=a_b^2\), and
\[
 4bu+p_b=b(4u+p)+(b-1)R.
\]
Thus the descendant is a selector state and its quotient is
\(T_b=bT+b-1\).  Since \(K=(T+1)/4\) in this channel,
\(K_b=(T_b+1)/4=bK\).  Applying~(3.4) gives
\[
 \ell_{0,b}=\frac{T_b+1}{4}=b\ell_0,
 \qquad
 \ell_{1,b}=bu-jT_b=b\ell_1-j(b-1),
\]
which proves~(7.3)--(7.6), including the intertwining identity~(7.4).
Each map is visibly injective for fixed \(b\).  Substitution into
~(7.1)--(7.3), with the source prime retained, proves the typed composition
law~(7.7).
\end{proof}

Reverse descent need not preserve divisor support: the factor \(b\) can
introduce divisors absent from a proposed parent.  Therefore~(7.2)--(7.7)
propagate witnesses but cannot create the first witness on an empty ray.

\subsection{Sharp verdict}

The strongest theorem supported by the audited data is
Theorem~\ref{thm:esksb-bridge} together with the exhaustive equivalence~(3.11):
the complete tagged Erd\H{o}s--Straus selector for a fixed prime
\(p\equiv1\pmod {12}\) is exactly the marked divisor-supported S6 terminal
slice, and the two integral relation matrices have the same row lattice.
The exact remaining obstruction is~(4.5).

Three stronger readings are false or unsupported.

\begin{enumerate}
\item The abstract group \(G(\ell)\), its determinant, or an unmarked
terminal triple does not determine selector support.  Every \(E\)-state has
the same trivial presented group, every \(M\)-state has group \(\ZZ/p\ZZ\),
and Examples~\ref{ex:esksb-tags} and~\ref{ex:esksb-ambient} exhibit the information lost by forgetting
the marker and divisor gate.

\item The row shear is not a geometric map from the reciprocal cubic
\(4xyz=p(xy+xz+yz)\) to the candidate complex threefold \(X\).  No map of
varieties or analytic spaces, no cycle map, no attaching map, and no
cohomological correspondence is defined by left multiplication of a
relation matrix.

\item The existence of the full S6 affine fibre cannot prove universal
Erd\H{o}s--Straus positivity.  Those fibres are always nonempty by~(5.2),
whereas selector nonemptiness is exactly the additional finite nonlinear
incidence in~(4.5).  The natural affine point is completely exhausted by
Theorem~\ref{thm:esksb-affine-gate} and contributes nothing to the unresolved
\(p\equiv1\pmod {24}\) half.
\end{enumerate}

Accordingly, the bridge is a genuine theorem and is stronger than a shared
determinant or a heuristic analogy.  It is also exact enough to prove its own
boundary: without a new argument forcing one of the bounded factor
congruences in~(4.5), it neither proves nor disproves the Erd\H{o}s--Straus
conjecture.
